
\documentclass{article}
\usepackage{colortbl}
\usepackage{makecell}
\usepackage{multirow}
\usepackage{supertabular}

\begin{document}

\newcounter{utterance}

\centering \large Interaction Transcript for game `eqbench', experiment `eq\_bench', episode 105 with reflexion{-}Qwen3.5{-}9B.
\vspace{24pt}

{ \footnotesize  \setcounter{utterance}{1}
\setlength{\tabcolsep}{0pt}
\begin{supertabular}{c@{$\;$}|p{.15\linewidth}@{}p{.15\linewidth}p{.15\linewidth}p{.15\linewidth}p{.15\linewidth}p{.15\linewidth}}
    \# & \multicolumn{2}{c}{Player} && \multicolumn{2}{c}{Game Master} \\
    \hline

    \theutterance \stepcounter{utterance}  
    & & & \multicolumn{4}{p{0.6\linewidth}}{
        \cellcolor[rgb]{0.9,0.9,0.9}{
            \makecell[{{p{\linewidth}}}]{
                \texttt{\tiny{[P1$\langle$GM]}}
                \texttt{You have attempted this same task before and did not succeed. These are the plans you wrote after those attempts.} \\
\texttt{Trial \#1: New plan: I failed because I assigned a high score to "Angry" based on Terry's defensive outburst, but I overlooked the underlying context of "survival" and "guilt" which points to deeper emotional turmoil. The phrase "Hit a nerve" suggests Cass struck a vulnerability, making "Insecure" or "Amused" (as a deflection) more likely than pure anger. Next time, I will analyze the dialogue for specific triggers: if a character is defending a plan as a matter of life or death while being mocked, the primary emotion is likely fear or insecurity masked as aggression. I will prioritize "Insecure" (score 8) and "Amused" (score 2, as a sarcastic deflection) over "Angry" (score 3), ensuring at least one score is high to reflect the intensity of the "nerve" being hit.} \\
\\ 
\texttt{Angry: 3} \\
\texttt{Insecure: 8} \\
\texttt{Amused: 2} \\
\texttt{Calm: 0} \\
\\ 
\\ 
\\ 
\texttt{Your task is to predict the likely emotional responses of a character in this dialogue:} \\
\\ 
\texttt{Terry: We're not going in blind, Cass. Not this time.} \\
\texttt{Cass: And who's gonna stop me? You?} \\
\texttt{Terry: If I have to. } \\
\texttt{Cass: Oh, right. The big, bad Terry Holloway. Always with a plan. } \\
\texttt{Terry: You're goddamn right. You know why? Because plans save lives.} \\
\texttt{Cass: Or maybe they just make you feel better about yourself. Like you're in control.} \\
\texttt{Terry: This isn't about control, Cass. It's about survival. } \\
\texttt{Cass: Survival? Or just your guilt?} \\
\texttt{Terry: That's enough.} \\
\texttt{Cass: Hit a nerve, did I? } \\
\texttt{{[}End dialogue{]}} \\
\\ 
\texttt{At the end of this dialogue, Terry would feel...} \\
\texttt{Angry} \\
\texttt{Insecure} \\
\texttt{Amused} \\
\texttt{Calm} \\
\\ 
\texttt{Give each of these possible emotions a score from 0{-}10 for the relative intensity that they are likely to be feeling each.} \\
\\ 
\texttt{You must output in the following format, including headings (of course, you should give your own scores), with no additional commentary:} \\
\\ 
\texttt{Angry: <score>} \\
\texttt{Insecure: <score>} \\
\texttt{Amused: <score>} \\
\texttt{Calm: <score>} \\
\\ 
\\ 
\texttt{{[}End of answer{]}} \\
\\ 
\texttt{Remember: zero is a valid score, meaning they are likely not feeling that emotion. You must score at least one emotion > 0.} \\
\\ 
\texttt{Your answer:} \\
            }
        }
    }
     \\ \\

    \theutterance \stepcounter{utterance}  
    & \multicolumn{4}{p{0.6\linewidth}}{
        \cellcolor[rgb]{0.561,0.737,0.561}{
            \makecell[{{p{\linewidth}}}]{
                \texttt{\tiny{[P1$\rangle$GM]}}
                \texttt{Angry: 4} \\
\texttt{Insecure: 9} \\
\texttt{Amused: 1} \\
\texttt{Calm: 0} \\
            }
        }
    }
    & & \\ \\

    \theutterance \stepcounter{utterance}  
    & & & \multicolumn{2}{p{0.3\linewidth}}{
        \cellcolor[rgb]{0.9,0.9,0.9}{
            \makecell[{{p{\linewidth}}}]{
                \texttt{\tiny{[GM$|$GM]}}
                \texttt{Angry: 4} \\
\texttt{Insecure: 9} \\
\texttt{Amused: 1} \\
\texttt{Calm: 0} \\
            }
        }
    }
    & & \\ \\

    \theutterance \stepcounter{utterance}  
    & & & \multicolumn{2}{p{0.3\linewidth}}{
        \cellcolor[rgb]{0.9,0.9,0.9}{
            \makecell[{{p{\linewidth}}}]{
                \texttt{\tiny{[GM$|$GM]}}
                \texttt{Angry: 8} \\
\texttt{Insecure: 0} \\
\texttt{Amused: 0} \\
\texttt{Calm: 0} \\
            }
        }
    }
    & & \\ \\

    \theutterance \stepcounter{utterance}  
    & & & \multicolumn{2}{p{0.3\linewidth}}{
        \cellcolor[rgb]{0.9,0.9,0.9}{
            \makecell[{{p{\linewidth}}}]{
                \texttt{\tiny{[GM$|$GM]}}
                \texttt{game\_result = LOSE} \\
            }
        }
    }
    & & \\ \\

    \theutterance \stepcounter{utterance}  
    & & \multicolumn{4}{p{0.6\linewidth}}{
        \cellcolor[rgb]{0.95,0.95,0.95}{
            \makecell[{{p{\linewidth}}}]{
                \texttt{\tiny{[P1$\leftrightarrow$P]}}
                \texttt{Post{-}episode Reflexion (not sent to GM):} \\
\texttt{New plan: I failed because I misinterpreted Terry's "That's enough" as a sign of insecurity rather than a command to end the conversation, and I underestimated the intensity of his frustration. Terry has been defending his plan as a matter of "survival" and "life," and Cass's mockery ("guilt," "feel better about yourself") directly attacks the core of his motivation. His final line is a sharp cutoff, indicating he has reached his limit of tolerance, not a moment of self{-}doubt. Next time, I will recognize that when a character defends a critical action against personal attacks and then abruptly shuts down the interlocutor, the dominant emotion is high{-}intensity anger/frustration. I will assign a high score to "Angry" (9) to reflect the defensive outburst and the final cutoff, while lowering "Insecure" (3) to represent the underlying vulnerability that was triggered but is currently suppressed by rage.} \\
\\ 
\texttt{Angry: 9} \\
\texttt{Insecure: 3} \\
\texttt{Amused: 0} \\
\texttt{Calm: 0} \\
            }
        }
    }
    & \\ \\

\end{supertabular}
}

\end{document}
